\begin{proof}[Proof of \cref{obj:thm:sharp-length-frontier-full-elbow}]
\begin{enumerate}
\item By \cref{obj:thm:honest-upper-full-elbow}, there is a finite constant \(C_0>0\), depending only on \((\alpha,c_f,C_f,L)\), such that
\[
I_n\in\mathfrak H_n
\qquad\text{for every integer }n>0.
\]
For every integer \(n\ge256\), this same constant satisfies
\[
\mathbb E_P\!\left[\lambda_{\Theta}(I_n)\right]
\le C_0\min\left\{1,\frac{n^{-1/3}}{\mu(P)}\right\}
\qquad\text{for every }P\in\mathcal M_n,
\]
and
\[
L_n(a)\le C_0\min\left\{1,\frac{n^{-1/3}}{a}\right\}
\qquad\text{for every }0<a\le\frac14.
\]
% lean: CausalSmith.Stat.TransportCaceRoughnuisanceLength.honest_upper_full_elbow

\item By \cref{obj:thm:honest-lower-full-elbow}, choose a positive lower comparison constant \(c_{\mathrm{lower}}\), depending only on \((\alpha,c_f,C_f,L)\), for which
\[
c_{\mathrm{lower}}\min\left\{1,\frac{n^{-1/3}}{a}\right\}
\le L_n(a)
\qquad
\text{whenever }n\ge256,\quad 0<a\le\frac14.
\]
% lean: CausalSmith.Stat.TransportCaceRoughnuisanceLength.honest_lower_full_elbow

\item Define the final lower comparison constant by
\[
c_0:=\min\left\{c_{\mathrm{lower}},\frac{C_0}{2}\right\}.
\]
Both entries in this minimum are positive. Consequently,
\[
0<c_0,
\qquad
c_0\le\frac{C_0}{2}<C_0.
\]
Thus \(c_0\) and \(C_0\) have the required dependence and strict ordering.
% lean: CausalSmith.Stat.TransportCaceRoughnuisanceLength.sharp_length_frontier_full_elbow

\item Fix an integer \(n\ge256\) and a real \(a\) with \(0<a\le1/4\). Since \(n>0\) and \(a>0\),
\[
0\le\min\left\{1,\frac{n^{-1/3}}{a}\right\}.
\]
Multiplying \(c_0\le c_{\mathrm{lower}}\) by this nonnegative factor and applying the lower bound gives
\[
c_0\min\left\{1,\frac{n^{-1/3}}{a}\right\}
\le
c_{\mathrm{lower}}\min\left\{1,\frac{n^{-1/3}}{a}\right\}
\le L_n(a).
\]
Together with the upper bound already obtained, this yields
\[
c_0\min\left\{1,\frac{n^{-1/3}}{a}\right\}
\le L_n(a)
\le C_0\min\left\{1,\frac{n^{-1/3}}{a}\right\}.
\]
The honesty assertion and the law-specific expected-length bound follow from \cref{obj:thm:honest-upper-full-elbow} with the same interval sequence and the same \(C_0\).
% lean: CausalSmith.Stat.TransportCaceRoughnuisanceLength.sharp_length_frontier_full_elbow
\end{enumerate}
\end{proof}

\begin{proof}[Proof of \cref{obj:thm:honest-upper-full-elbow}]
\begin{enumerate}
\item
Let \(C_{\mathrm{mse}}\) be the mean-square envelope in
\cref{obj:lem:marked-cubic-mse}. Its displayed expression is finite and strictly positive under the stated conditions. For every positive integer \(n\), define the calibration radius by
\[
r_n:=\sqrt{\frac{2C_{\mathrm{mse}}}{\alpha}}\,n^{-1/3}\ge0.
\]
Thus the acceptance set in \cref{obj:synth_10} is
\[
\mathcal A_n(\omega)
=
\left\{\vartheta\in[-1,1]:
\left|\widehat T_Y(\omega)-\vartheta\widehat T_D(\omega)\right|
\le2r_n\right\},
\]
where the two statistics are those of \cref{obj:def:cubic-estimator}.

The statistics are measurable. Moreover,
\[
\mathcal A_n(\omega)\ne\varnothing
\quad\Longleftrightarrow\quad
|\widehat T_Y(\omega)|
\le|\widehat T_D(\omega)|+2r_n.
\]
Indeed, the values of \(\vartheta\widehat T_D(\omega)\) for
\(\vartheta\in[-1,1]\) form
\([ -|\widehat T_D(\omega)|,|\widehat T_D(\omega)|]\), whose distance from
\(\widehat T_Y(\omega)\) is at most \(2r_n\) precisely under the displayed inequality. Consequently, both the nonemptiness event and the graph of \(\mathcal A_n\) are measurable. For \(n<n_0\), the graph of \(I_n\) is the product of the sample space with \(\Theta=[-1,1]\). For \(n\ge n_0\), its graph combines the graph of \(\mathcal A_n\) on the nonemptiness event with the graph of \(\{0\}\) on its complement. Hence
\[
\{(\omega,\vartheta):\vartheta\in I_n(\omega)\}
\quad\text{is measurable}.
\]
Under every \(P\in\mathcal M_n\), the joint sampling law is a probability law, and the squared estimation errors are integrable for both responses. Below the threshold the errors are constant; above it their integrability follows from the square-integrable centered estimator and conditional-mean error.
% lean: CausalSmith.Stat.TransportCaceRoughnuisanceLength.cubicEstimator_error_sq_integrable
% lean: CausalSmith.Stat.TransportCaceRoughnuisanceLength.scoreInterval_measurableSet_graph

\item
We first check connectedness. Fix a realization \(\omega\), and take
\(\vartheta_1,\vartheta_2\in\mathcal A_n(\omega)\) with
\(\vartheta_1\le\vartheta_2\). For every
\(\vartheta\in[\vartheta_1,\vartheta_2]\), if
\(\widehat T_D(\omega)\ge0\), then
\[
-2r_n
\le \widehat T_Y(\omega)-\vartheta_2\widehat T_D(\omega)
\le \widehat T_Y(\omega)-\vartheta\widehat T_D(\omega)
\le \widehat T_Y(\omega)-\vartheta_1\widehat T_D(\omega)
\le2r_n.
\]
If \(\widehat T_D(\omega)<0\), the corresponding inequalities are
\[
-2r_n
\le \widehat T_Y(\omega)-\vartheta_1\widehat T_D(\omega)
\le \widehat T_Y(\omega)-\vartheta\widehat T_D(\omega)
\le \widehat T_Y(\omega)-\vartheta_2\widehat T_D(\omega)
\le2r_n.
\]
Also \(\vartheta\in[-1,1]\), so \(\vartheta\in\mathcal A_n(\omega)\).
The full-domain and singleton branches in \cref{obj:def:score-interval}
are connected subsets of \(\Theta\) as well. Thus, for every realization,
\[
I_n(\omega)\subseteq\Theta,
\qquad I_n(\omega)\ \text{is connected}.
\]
% lean: CausalSmith.Stat.TransportCaceRoughnuisanceLength.scoreInterval_subset_ordConnected

For coverage, fix \(P\in\mathcal M_n\).
\Cref{obj:lem:transported-cace-identification} supplies
\[
\mu(P)=T_D(P)>0,\qquad
\theta(P)=\frac{T_Y(P)}{\mu(P)}\in[-1,1],
\qquad
T_Y(P)=\theta(P)\mu(P).
\]
If \(n<n_0\), the interval equals \(\Theta\), and therefore
\[
P\{\theta(P)\in I_n\}=1\ge1-\alpha.
\]
Suppose now that \(n\ge n_0\). For each \(A\in\{Y,D\}\), define the error event
\[
\mathcal E_{A,n}
:=
\{\omega:|\widehat T_A(\omega)-T_A(P)|>r_n\}.
\]
The radius is strictly positive and satisfies
\[
r_n^2=\frac{2C_{\mathrm{mse}}}{\alpha}\,n^{-2/3}.
\]
Applying the second-moment tail inequality and
\cref{obj:lem:marked-cubic-mse} gives
\[
P(\mathcal E_{A,n})
\le
\frac{\mathbb E_P[(\widehat T_A-T_A(P))^2]}{r_n^2}
\le
\frac{C_{\mathrm{mse}}n^{-2/3}}
     {(2C_{\mathrm{mse}}/\alpha)n^{-2/3}}
=\frac{\alpha}{2}.
\]
% lean: CausalSmith.Stat.TransportCaceRoughnuisanceLength.cubicEstimator_error_tail_le

Define the simultaneous good event by
\[
\mathcal E_{\mathrm{cov},n}
:=(\mathcal E_{Y,n}\cup\mathcal E_{D,n})^c.
\]
On this event, identification and the triangle inequality yield
\[
\begin{aligned}
|\widehat T_Y-\theta(P)\widehat T_D|
&=
|(\widehat T_Y-T_Y(P))
 +\theta(P)(\mu(P)-\widehat T_D)|\\
&\le|\widehat T_Y-T_Y(P)|
  +|\theta(P)|\,|\widehat T_D-\mu(P)|\\
&\le2r_n.
\end{aligned}
\]
Hence \(\theta(P)\in\mathcal A_n\), so the acceptance set is nonempty and
\(\theta(P)\in I_n\). The union bound therefore gives
\[
P\{\theta(P)\in I_n\}
\ge P(\mathcal E_{\mathrm{cov},n})
\ge1-P(\mathcal E_{Y,n})-P(\mathcal E_{D,n})
\ge1-\alpha.
\]
Together with the measurable graph and connectedness established above,
this proves \(I_n\in\mathfrak H_n\) for every positive integer \(n\), by
\cref{obj:def:honest-intervals}.
% lean: CausalSmith.Stat.TransportCaceRoughnuisanceLength.scoreInterval_coverage_of_square_integrable

\item
Choose
\[
C_0:=
\max\left\{
2,\,
8\sqrt{\frac{2C_{\mathrm{mse}}}{\alpha}}
+8C_{\mathrm{mse}}
\right\}.
\]
This constant is finite, depends only on
\((\alpha,c_f,C_f,L)\), and satisfies
\[
C_0\ge2>0.
\]
% lean: CausalSmith.Stat.TransportCaceRoughnuisanceLength.honest_upper_full_elbow

\item
Fix \(n\ge n_0\) and \(P\in\mathcal M_n\). Define the bad-slope event by
\[
\mathcal E_{\mathrm{slope},n,P}
:=
\{\omega:|\widehat T_D(\omega)-\mu(P)|>\mu(P)/2\}.
\]
Since \(\mu(P)>0\), the second-moment tail inequality and
\cref{obj:lem:marked-cubic-mse} imply
\[
\begin{aligned}
P(\mathcal E_{\mathrm{slope},n,P})
&\le
\frac{\mathbb E_P[(\widehat T_D-\mu(P))^2]}
     {(\mu(P)/2)^2}\\
&\le
\frac{4C_{\mathrm{mse}}n^{-2/3}}{\mu(P)^2}.
\end{aligned}
\]
% lean: CausalSmith.Stat.TransportCaceRoughnuisanceLength.cubicEstimator_bad_slope_probability_le

On the complementary event, the reverse triangle inequality gives
\[
|\widehat T_D|
\ge \mu(P)-|\widehat T_D-\mu(P)|
\ge\frac{\mu(P)}2>0.
\]
For a nonzero slope, the unrestricted affine-score acceptance interval is
\[
\left\{\vartheta\in\mathbb R:
|\widehat T_Y-\vartheta\widehat T_D|\le2r_n\right\}
=
\left[
\frac{\widehat T_Y}{\widehat T_D}
-\frac{2r_n}{|\widehat T_D|},
\,
\frac{\widehat T_Y}{\widehat T_D}
+\frac{2r_n}{|\widehat T_D|}
\right].
\]
Its length is \(4r_n/|\widehat T_D|\). Intersecting with \(\Theta\) bounds
the length by both this quantity and \(2\). Moreover,
\[
\lambda_\Theta(I_n)=\lambda_\Theta(\mathcal A_n):
\]
the two sets agree when \(\mathcal A_n\) is nonempty, and otherwise both
the empty acceptance set and the singleton fallback have length zero.
Consequently, on the complementary event,
\[
\lambda_\Theta(I_n)
\le
\min\left\{2,\frac{4r_n}{|\widehat T_D|}\right\}
\le
\min\left\{2,\frac{8r_n}{\mu(P)}\right\}.
\]
% lean: CausalSmith.Stat.TransportCaceRoughnuisanceLength.restrictedLength_scoreInterval_le_on_good_slope

Define the nonnegative length majorant
\[
\ell_{\mathrm{good},n,P}:=\frac{8r_n}{\mu(P)}\ge0.
\]
Restricted length always lies between zero and the length of \(\Theta\).
Thus, on the bad-slope event it is at most \(2\), while on its complement
the preceding bound applies. For every realization,
\[
0\le\lambda_\Theta(I_n)
\le
\ell_{\mathrm{good},n,P}
+2\mathbf1_{\mathcal E_{\mathrm{slope},n,P}}.
\]
Integrating this inequality gives
\[
\begin{aligned}
\mathbb E_P[\lambda_\Theta(I_n)]
&\le
\ell_{\mathrm{good},n,P}
+2P(\mathcal E_{\mathrm{slope},n,P})\\
&\le
8\sqrt{\frac{2C_{\mathrm{mse}}}{\alpha}}\,
\frac{n^{-1/3}}{\mu(P)}
+
8C_{\mathrm{mse}}\frac{n^{-2/3}}{\mu(P)^2}.
\end{aligned}
\]
There is also the deterministic cap
\[
\mathbb E_P[\lambda_\Theta(I_n)]\le2.
\]
% lean: CausalSmith.Stat.TransportCaceRoughnuisanceLength.scoreInterval_expectedLength_le_of_square_integrable

Set
\[
\xi_{n,P}:=\frac{n^{-1/3}}{\mu(P)}>0,
\qquad
\xi_{n,P}^2=\frac{n^{-2/3}}{\mu(P)^2}.
\]
If \(\xi_{n,P}\le1\), then \(\xi_{n,P}^2\le\xi_{n,P}\), and hence
\[
\begin{aligned}
\mathbb E_P[\lambda_\Theta(I_n)]
&\le
\left(
8\sqrt{\frac{2C_{\mathrm{mse}}}{\alpha}}
+8C_{\mathrm{mse}}
\right)\xi_{n,P}\\
&\le C_0\xi_{n,P}
=C_0\min\{1,\xi_{n,P}\}.
\end{aligned}
\]
If \(\xi_{n,P}>1\), the cap gives
\[
\mathbb E_P[\lambda_\Theta(I_n)]
\le2\le C_0
=C_0\min\{1,\xi_{n,P}\}.
\]
This proves the asserted law-specific expected-length bound.
% lean: CausalSmith.Stat.TransportCaceRoughnuisanceLength.upper_length_two_regime_bound

\item
Fix \(n\ge n_0\) and \(0<a\le1/4\).
The proposed slice bound is nonnegative:
\[
0\le C_0\min\left\{1,\frac{n^{-1/3}}a\right\}.
\]
For every \(P\in\mathcal M_n(a)\),
\cref{obj:def:strength-slice} supplies \(P\in\mathcal M_n\) and
\(a\le\mu(P)\). Since the numerator is nonnegative and \(a>0\),
\[
\frac{n^{-1/3}}{\mu(P)}\le\frac{n^{-1/3}}a.
\]
The law-specific bound therefore yields
\[
\mathbb E_P[\lambda_\Theta(I_n)]
\le
C_0\min\left\{1,\frac{n^{-1/3}}{\mu(P)}\right\}
\le
C_0\min\left\{1,\frac{n^{-1/3}}a\right\}.
\]
If the slice is nonempty, taking its supremum preserves this bound.
For an empty slice, the supremum is taken to be zero, and the same
conclusion follows from the nonnegativity displayed above. In either case,
\[
\sup_{P\in\mathcal M_n(a)}
\mathbb E_P[\lambda_\Theta(I_n)]
\le
C_0\min\left\{1,\frac{n^{-1/3}}a\right\}.
\]
% lean: CausalSmith.Stat.TransportCaceRoughnuisanceLength.honest_upper_full_elbow

\item
To pass to the infimum, observe that every
\(C_n\in\mathfrak H_n\) and every \(P\in\mathcal M_n(a)\) satisfy
\[
0\le\lambda_\Theta(C_n(\omega))\le2,
\qquad
0\le\mathbb E_P[\lambda_\Theta(C_n)]\le2.
\]
These inequalities follow from nonnegativity of Lebesgue length,
restriction to \(\Theta=[-1,1]\), and the probability sampling law.
For a nonempty slice, the expected lengths consequently have a finite
supremum satisfying
\[
0\le
\sup_{P\in\mathcal M_n(a)}
\mathbb E_P[\lambda_\Theta(C_n)]
\le2.
\]
For an empty slice this supremum equals zero. Thus the collection of
suprema over honest procedures is bounded below by zero. It is also
nonempty, since \(I_n\in\mathfrak H_n\). Using \(I_n\) as a competitor in
\cref{obj:def:length-frontier} gives
\[
\begin{aligned}
L_n(a)
&=
\inf_{C_n\in\mathfrak H_n}
\sup_{P\in\mathcal M_n(a)}
\mathbb E_P[\lambda_\Theta(C_n)]\\
&\le
\sup_{P\in\mathcal M_n(a)}
\mathbb E_P[\lambda_\Theta(I_n)]\\
&\le
C_0\min\left\{1,\frac{n^{-1/3}}a\right\}.
\end{aligned}
\]
% lean: CausalSmith.Stat.TransportCaceRoughnuisanceLength.honest_upper_full_elbow
\end{enumerate}
\end{proof}

\begin{proof}[Proof of \cref{obj:thm:honest-lower-full-elbow}]
\begin{enumerate}
\item Choose the amplitude supplied by \cref{obj:lem:legal-iv-mixture-full-elbow}, and define
\[
0<c_{\star}<1,
\qquad
c_0=(1-\alpha)2^{3/4}c_{\star}^2>0.
\]
The amplitude is chosen uniformly in \(n\) and \(a\), so \(c_0\) depends on the prescribed constants \(\alpha,c_f,C_f,L\).
% lean: CausalSmith.Stat.TransportCaceRoughnuisanceLength.legal_iv_mixture_full_elbow

\item Fix \(n\ge256\) and \(0<a\le1/4\). Use the center, components, and observed mixtures of \cref{obj:lem:legal-iv-mixture-full-elbow}, with
\[
K_{\mathrm L}=\lceil n^{4/3}\rceil,
\qquad
\bar a=\max\{a,n^{-1/3}\},
\qquad
r_{\mathrm{sep}}=\frac{J_n}{\bar a}.
\]
Since the negative power is decreasing on the positive half-line,
\[
0<n^{-1/3}\le64^{-1/3}=\frac14,
\qquad
0<\bar a\le\frac14.
\]
The separation bound in that same result gives
\[
0<2^{3/4}c_{\star}^2n^{-1/3}\le J_n,
\qquad
r_{\mathrm{sep}}>0.
\]
For every sign vector \(\lambda\in\{-1,1\}^{K_{\mathrm L}}\), the component at \(\tau=-1\) belongs to \(\mathcal M_n(a)\subseteq\mathcal M_n\). Applying \cref{obj:lem:transported-cace-identification} to this component yields
\[
r_{\mathrm{sep}}
=\theta(Q_{\lambda,-1})
\in[-1,1].
\]
Consequently,
\[
0<r_{\mathrm{sep}}\le1,
\qquad
[-r_{\mathrm{sep}},r_{\mathrm{sep}}]\subseteq\Theta.
\]
The sampling conditions ensure that the center and component observed laws are probability laws.
% lean: CausalSmith.Stat.TransportCaceRoughnuisanceLength.actualStrength_bounds
% lean: CausalSmith.Stat.TransportCaceRoughnuisanceLength.transported_cace_identification

\item The strength floor implements the capped inverse-strength rate. Indeed, if \(a\le n^{-1/3}\), then
\[
\bar a=n^{-1/3},
\qquad
\frac{n^{-1/3}}{\bar a}=1
=\min\left\{1,\frac{n^{-1/3}}a\right\}.
\]
If \(a>n^{-1/3}\), then
\[
\bar a=a,
\qquad
\frac{n^{-1/3}}{\bar a}
=\frac{n^{-1/3}}a
=\min\left\{1,\frac{n^{-1/3}}a\right\}.
\]
Dividing the separation lower bound by \(\bar a>0\) and multiplying by \(1-\alpha>0\) therefore gives
\[
c_0\min\left\{1,\frac{n^{-1/3}}a\right\}
=(1-\alpha)2^{3/4}c_{\star}^2
  \frac{n^{-1/3}}{\bar a}
\le(1-\alpha)r_{\mathrm{sep}}.
\]
% lean: CausalSmith.Stat.TransportCaceRoughnuisanceLength.actualStrength_rate_identity

\item Let \(C_n\) be any procedure satisfying the theorem's hypotheses. For every \(t\in\mathbb R\), define the inclusion event
\[
E_t=\{\omega:t\in C_n(\omega)\}.
\]
These events are measurable by graph measurability. For \(t\in[-r_{\mathrm{sep}},r_{\mathrm{sep}}]\), define
\[
\tau_t=-\frac{t}{r_{\mathrm{sep}}}.
\]
The inequalities \(-r_{\mathrm{sep}}\le t\le r_{\mathrm{sep}}\), together with \(r_{\mathrm{sep}}>0\), give
\[
-1\le\tau_t\le1.
\]
By \cref{obj:lem:legal-iv-mixture-full-elbow}, every component indexed by \(\tau_t\) belongs to \(\mathcal M_n(a)\) and satisfies
\[
\theta(Q_{\lambda,\tau_t})
=-\frac{\tau_tJ_n}{\bar a}
=-\tau_t r_{\mathrm{sep}}
=t.
\]
Uniform honesty thus implies, for every sign vector,
\[
\left(
Q_{\lambda,\tau_t,\mathrm S}^{\otimes n}
\otimes Q_{\lambda,\tau_t,\mathrm T}^{\otimes n}
\right)(E_t)\ge1-\alpha.
\]
Averaging these inequalities over the finite sign family gives
\[
M_{\tau_t}(E_t)
=
2^{-K_{\mathrm L}}
\sum_{\lambda\in\{-1,1\}^{K_{\mathrm L}}}
\left(
Q_{\lambda,\tau_t,\mathrm S}^{\otimes n}
\otimes Q_{\lambda,\tau_t,\mathrm T}^{\otimes n}
\right)(E_t)
\ge1-\alpha.
\]
The normalized mixture is a probability law. The total-variation event inequality and the bound in \cref{obj:lem:legal-iv-mixture-full-elbow} yield
\[
M_{\tau_t}(E_t)-P_{\star}^{(n,n)}(E_t)
\le
\operatorname{TV}(M_{\tau_t},P_{\star}^{(n,n)})
\le\frac{1-\alpha}{2}.
\]
Hence, throughout the symmetric interval,
\[
P_{\star}^{(n,n)}(t\in C_n)
\ge\frac{1-\alpha}{2},
\qquad
t\in[-r_{\mathrm{sep}},r_{\mathrm{sep}}].
\]
% lean: CausalSmith.Stat.TransportCaceRoughnuisanceLength.uniform_data_mixture_coverage
% lean: Causalean.Stat.measureReal_sub_le_tvDist
% lean: CausalSmith.Stat.TransportCaceRoughnuisanceLength.honest_lower_full_elbow

\item Define the center's inclusion-probability function by
\[
p_{\star,n}(t)=P_{\star}^{(n,n)}(t\in C_n),
\qquad t\in\mathbb R.
\]
Graph measurability makes this function measurable, and
\[
0\le p_{\star,n}(t)\le1.
\]
It is therefore integrable over \(\Theta\). Applying \cref{obj:lem:random-set-length-identity} and using the interval inclusion established above gives
\[
\begin{aligned}
(1-\alpha)r_{\mathrm{sep}}
&=\int_{-r_{\mathrm{sep}}}^{r_{\mathrm{sep}}}
       \frac{1-\alpha}{2}\,dt\\
&\le\int_{-r_{\mathrm{sep}}}^{r_{\mathrm{sep}}}
       p_{\star,n}(t)\,dt\\
&\le\int_{\Theta}p_{\star,n}(t)\,dt\\
&=\mathbb E_{P_{\star}}
       \!\left[\lambda_{\Theta}(C_n)\right].
\end{aligned}
\]
The first inequality uses the pointwise inclusion bound; the second uses nonnegativity and
\([-r_{\mathrm{sep}},r_{\mathrm{sep}}]\subseteq\Theta\).
% lean: CausalSmith.Stat.TransportCaceRoughnuisanceLength.expectedLength_lower_of_inclusion

\item The range of expected lengths on the strength slice is bounded above. Pointwise,
\[
0\le\lambda_{\Theta}(C_n(\omega))
\le\operatorname{Leb}(\Theta)=2.
\]
For an integrable length, integration under an observed probability law gives
\[
\mathbb E_P\!\left[\lambda_{\Theta}(C_n)\right]
\le\int 2\,dP^{(n,n)}=2,
\qquad P\in\mathcal M_n(a).
\]
Under the convention assigning zero to a nonintegrable real integral, the complementary case also gives
\[
\mathbb E_P\!\left[\lambda_{\Theta}(C_n)\right]=0\le2.
\]
The measurable bounded lengths considered here are integrable. Since \(P_{\star}\in\mathcal M_n(a)\), the slice is nonempty, and comparison with its supremum now yields
\[
\begin{aligned}
c_0\min\left\{1,\frac{n^{-1/3}}a\right\}
&\le(1-\alpha)r_{\mathrm{sep}}\\
&\le\mathbb E_{P_{\star}}
       \!\left[\lambda_{\Theta}(C_n)\right]\\
&\le\sup_{P\in\mathcal M_n(a)}
       \mathbb E_P\!\left[\lambda_{\Theta}(C_n)\right]
\le2.
\end{aligned}
\]
This proves the asserted procedure-wise lower bound.
% lean: CausalSmith.Stat.TransportCaceRoughnuisanceLength.expectedLength_le_two
% lean: CausalSmith.Stat.TransportCaceRoughnuisanceLength.honest_lower_full_elbow

\item The honest interval class in \cref{obj:def:honest-intervals} is nonempty. Define the constant procedure by
\[
C_{n,\Theta}(\omega)=\Theta
\qquad\text{for every sample realization }\omega.
\]
It has a measurable graph and connected values. By \cref{obj:lem:transported-cace-identification},
\[
\theta(P)\in\Theta,
\qquad
P^{(n,n)}\bigl(\theta(P)\in C_{n,\Theta}\bigr)
=1\ge1-\alpha,
\qquad P\in\mathcal M_n.
\]
Thus
\[
C_{n,\Theta}\in\mathfrak H_n,
\qquad
\mathfrak H_n\ne\varnothing.
\]
Every member of this class satisfies the procedure-wise bound already established. Taking the infimum prescribed by \cref{obj:def:length-frontier} gives
\[
L_n(a)
=
\inf_{C_n\in\mathfrak H_n}
\sup_{P\in\mathcal M_n(a)}
\mathbb E_P\!\left[\lambda_{\Theta}(C_n)\right]
\ge
c_0\min\left\{1,\frac{n^{-1/3}}a\right\}.
\]
% lean: CausalSmith.Stat.TransportCaceRoughnuisanceLength.parameterSpace_honest
% lean: CausalSmith.Stat.TransportCaceRoughnuisanceLength.honest_lower_full_elbow
\end{enumerate}
\end{proof}
